\item El núcleo de un átomo se puede modelar como constituido por protones y neutrones cercanamente empaquetados. Cada partícula tiene una masa de $1.67 \times 10^{-27} \text{ kg}$ y un radio del orden de $10^{-15} \text{ m}$. (a) Utilice este modelo y los datos proporcionados para estimar la densidad del núcleo de un átomo. (b) Compare su resultado con la densidad de un material, por ejemplo, el hierro. ¿Qué sugieren su resultado y la comparación respecto a la estructura de la materia?
